\section{ALC}
\label{section:alc}

\begin{frame}{Automatic Level Control}
    \begin{columns}
        \begin{column}{0.48\linewidth}
            \DARCimage{\linewidth}{915include}
        \end{column}
        \begin{column}{0.48\linewidth}
            \begin{itemize}
                \item \emph{Automatic-Level-Control} (ALC) regelt die Aussteuerung der Senderendstufe
                \item Reduziert bei Übersteuerung die Amplitude im Sendezweig
                \item Nicht verwechseln mit der AGC (Automatic-Gain-Control) im Empfängerzweig
                \item ALC-Anzeige ist deshalb oft ein Balken
            \end{itemize}
        \end{column}
    \end{columns}
\end{frame}

\begin{frame}{Funktionsweise ALC}
    \begin{center}
        \DARCimage{0.75\linewidth}{914include}
    \end{center}
    \begin{itemize}
        \item Erfassung der Ausgangsleistung und Vergleich mit vorgegebenen maximalen Wert
        \item Bei Überschreiten wird eine Regelspannung an vorgelagerte HF-Verstärkerstufe gegeben
        \item Amplitude des Sendesignals wird reduziert
    \end{itemize}
\end{frame}

\begin{frame}{Anmerkungen zur ALC}
    \begin{itemize}
        \item Wenn ALC-Anzeige anspricht, ist Sender durch zu starkes NF-Signal
        übersteuert
        \item Bei SSB ist ein geringfügiges Übersteuern erwünscht, da eine
        kompression vorgenommen wird
        \item Bei Digimodes mit fester Amplitude, wie FT8, ist das Eingreifen
        der ALC erstmal nicht schlimm, könnte eine ansprechende ALC ein Indiz
        sein, dass das NF-Signal übersteuert ist
        \item NF-Signal überprüfen, dass es sauber und nicht übersteuert ist
        \item Bei Digimodes wie PSK31 ist die durch die ALC versuchachte
        \enquote{AM} störend
    \end{itemize}
\end{frame}

\begin{frame}
    \only<1>{
        \begin{QQuestion}{EF305}{Was bewirkt die ALC (Automatic Level Control) bei zu starkem NF-Signal in einem Transceiver?}{Sie reduziert die Amplitude des Signals im Sendezweig vor dem Leistungsverstärker.}
            {Sie erhöht die Amplitude des Signals im Sendezweig vor dem Leistungsverstärker.}
            {Sie reduziert die Verstärkung von Verstärkerstufen im Empfangsteil.}
            {Sie erhöht die Verstärkung von Verstärkerstufen im Empfangsteil.}
        \end{QQuestion}

    }
    \only<2>{
        \begin{QQuestion}{EF305}{Was bewirkt die ALC (Automatic Level Control) bei zu starkem NF-Signal in einem Transceiver?}{\textbf{\textcolor{DARCgreen}{Sie reduziert die Amplitude des Signals im Sendezweig vor dem Leistungsverstärker.}}}
            {Sie erhöht die Amplitude des Signals im Sendezweig vor dem Leistungsverstärker.}
            {Sie reduziert die Verstärkung von Verstärkerstufen im Empfangsteil.}
            {Sie erhöht die Verstärkung von Verstärkerstufen im Empfangsteil.}
        \end{QQuestion}

    }
\end{frame}%ENDCONTENT